%────────────────────── Capítulo 5 — Conclusiones ──────────────────────
% Resumen con los puntos fuertes y débiles de la vulnerabilidad.
\chapter{Conclusiones}
\label{cap:conclusiones}

El análisis de consumo por correlación extrae la clave de un AES-128 sin
atacar el algoritmo: aprovecha que una implementación concreta consume de forma
distinta según los datos que procesa. Este capítulo resume qué hace fuerte a
esa vulnerabilidad, qué necesita el atacante para explotarla y qué queda
abierto. En este trabajo se recorre el camino completo con trazas medidas: en
un ChipWhisperer-Nano, con una clave de laboratorio y el AES de NewAE sin
protecciones, el CPA recupera los 16 bytes con \num{1000} trazas, y el 90~\% de
las ordenaciones aleatorias de esa campaña lo consigue con \num{112}
(\cref{cap:resultados}). La contramedida de la fase 6 muestra el otro lado: el
barajado encarece el ataque pero no lo impide, y los retardos aleatorios lo
tumban tal como está escrito, aunque un atacante dispuesto a preprocesar las
trazas recupera parte de la clave.

\section{Puntos fuertes del ataque}
\label{sec:puntos-fuertes}

\begin{itemize}
  \item \textbf{Divide la clave.} El valor intermedio
    $\mathrm{SBox}(p_j \oplus k_j)$ depende de un solo byte de clave, así que
    la búsqueda de $2^{128}$ claves se reduce a $16 \times 256 = 4096$
    hipótesis que se comprueban por separado (\cref{sec:aes}). La seguridad
    matemática de AES no interviene.
  \item \textbf{Solo necesita textos claros y trazas.} El atacante no lee la
    memoria ni modifica el software, y el dispositivo cifra con normalidad
    mientras se mide. Basta conocer los textos claros, o los criptogramas si
    se ataca la última ronda.
  \item \textbf{Aprovecha toda la fuga del byte.} Frente al DPA de un bit, el
    CPA usa los ocho bits del valor intermedio y no obliga a adivinar qué bit
    filtra más \cite{brier2004cpa}. En la simulación, con 2000 trazas y
    $\mathrm{SNR} = 0{,}1$, el CPA recupera los 16 bytes y el DPA, entre 7 y 15
    según el bit (\cref{sec:validacion-simulada}; datos simulados).
  \item \textbf{Escala de forma predecible con el ruido.} Como el ruido que no
    depende de los datos se promedia, el número de trazas crece
    aproximadamente como $1/\mathrm{SNR}$ (\cref{tab:sim-ruido}; datos
    simulados). Más ruido encarece el ataque, pero no lo impide.
  \item \textbf{El código de tiempo constante no lo frena.} Un AES que ejecuta
    siempre las mismas instrucciones produce trazas alineadas, que es justo lo
    que el CPA necesita (\cref{sec:tiempo-constante}).
  \item \textbf{Equipo asequible.} Existen plataformas de bajo coste pensadas
    para enseñar estos ataques, como el ChipWhisperer-Nano que se usa en este
    trabajo \cite{cw-nano-docs}, y el cálculo cabe en un portátil.
  \item \textbf{Es general.} El mismo procedimiento sirve para cualquier
    cifrado en el que un paso no lineal mezcle el dato conocido con una parte
    pequeña de la clave \cite{mangard2007power}.
\end{itemize}

\section{Puntos débiles y requisitos}
\label{sec:puntos-debiles}

\begin{itemize}
  \item \textbf{Acceso físico.} Hay que medir el consumo del chip, lo que en un
    dispositivo real supone intervenir en su alimentación o acercar una sonda
    electromagnética. Es un ataque a dispositivos que el atacante tiene en la
    mano, no a distancia.
  \item \textbf{Muchas ejecuciones con la misma clave.} El CPA necesita cientos
    o miles de cifrados con la misma clave y textos conocidos. Un protocolo que
    renueve la clave antes de ese umbral lo deja sin datos
    (\cref{sec:rekeying}).
  \item \textbf{Sincronización.} El CPA compara las trazas instante a instante.
    En el laboratorio el propio firmware da la señal de disparo; un atacante
    real tiene que encontrar el cifrado en la traza y alinear las trazas por su
    cuenta, y los retardos aleatorios o el \emph{shuffling} empeoran ese
    problema a propósito (\cref{sec:shuffling}).
  \item \textbf{Un modelo de fuga adecuado.} El peso de Hamming es una
    aproximación. Si el dispositivo filtra sobre todo la distancia de Hamming
    entre valores consecutivos, el modelo pierde correlación y hacen falta más
    trazas, o un modelo distinto \cite{brier2004cpa}.
  \item \textbf{Las contramedidas funcionan.} Un enmascaramiento correcto anula
    el CPA de primer orden sea cual sea el número de trazas, y obliga a ataques
    de orden superior, bastante más caros (\cref{sec:masking}).
  \item \textbf{Condiciones favorables en este trabajo.} La clave y el firmware
    los elige el grupo, la señal de disparo la da el firmware y el capturador
    comparte reloj con el microcontrolador. Son las condiciones de un
    laboratorio de enseñanza, no las de un ataque a un producto, y los números
    que se obtengan no se pueden trasladar sin más a otro dispositivo.
\end{itemize}
En la placa, el ruido de la medida resultó ser del orden de un escalón del
convertidor y las trazas salieron alineadas al ciclo sin procesado, así que
ninguno de los problemas habituales de una captura real ha estorbado
(\cref{sec:calidad-captura}). Las \num{112} trazas para el 90~\% de los ensayos
están dentro de lo que predecía el modelo simulado para la relación
señal-ruido medida, del lado bueno: la fuga real aparece en más muestras que en
la simulación, y eso también ayuda al atacante.

\section{Trabajo futuro}
\label{sec:trabajo-futuro}

\begin{itemize}
  \item \textbf{Completar la contramedida.} Las variantes medidas usan un
    generador xorshift32 con semilla fija, que un atacante que conozca el
    código puede predecir (\cref{sec:contramedidas-medidas}). Haría falta un
    generador realmente aleatorio en el dispositivo, y un enmascaramiento, que
    es la única contramedida de las estudiadas que anula el CPA de primer
    orden. Del lado del atacante, queda por escribir un ataque completo de
    realineado contra los retardos, que en este trabajo se ha compensado solo
    con ventanas de integración.
  \item \textbf{Otros modelos y objetivos.} Atacar la última ronda sobre trazas
    reales y comparar el peso de Hamming con la distancia de Hamming.
  \item \textbf{Ataques de orden superior.} Repetir el experimento contra una
    implementación enmascarada y probar un DPA de segundo orden
    \cite{messerges2000second}.
  \item \textbf{Ataques de plantillas.} Caracterizar la fuga con un dispositivo
    controlado para atacar otro con muy pocas trazas
    \cite[sec.~II.B.4]{mishra2025modern}.
  \item \textbf{Un montaje propio.} Repetir la medida con un microcontrolador
    comercial y un \emph{shunt} en su alimentación, sin una plataforma pensada
    para el ataque, como el ATmega328P en placa de prototipos que se planteó al
    principio del trabajo.
  \item \textbf{Canal electromagnético.} Medir la radiación del chip con una
    sonda de campo cercano en lugar del consumo en la alimentación
    \cite[sec.~II.B]{mishra2025modern}, lo que puede evitar tener que
    intervenir en la placa.
\end{itemize}
